% !TEX program = pdflatex
\documentclass[11pt]{article}
\usepackage[margin=1in]{geometry}
\usepackage{amsmath,amssymb,amsthm}
\usepackage{hyperref}
\newtheorem{theorem}{Theorem}

\title{A Short Note on Gaussian Sums}
\author{A. Author}
\date{}

\begin{document}
\maketitle

\begin{abstract}
We study a mixture of Gaussian densities and show that the bound is tight
for all $n$. This note is a sample document for texview.
\end{abstract}

\section{Introduction}
\label{sec:intro}
A Gaussian sum approximates a density $p(x)$ by a weighted sum of normal
densities \cite{alspach1972}. The weights are positive and sum to one.

\section{Main result}
Let $w_i \ge 0$ with $\sum_i w_i = 1$. Then
\begin{equation}
  p(x) \approx \sum_{i=1}^{n} w_i \, \mathcal{N}(x; \mu_i, \Sigma_i).
  \label{eq:sum}
\end{equation}

\begin{theorem}
The error of \eqref{eq:sum} goes to zero as $n \to \infty$.
\end{theorem}

\begin{proof}
See Section~\ref{sec:intro} and the references therein.
\end{proof}

\section{Conclusion}
The bound is tight for all $n$, and the method is simple to use.

\bibliographystyle{plain}
\bibliography{paper}
\end{document}
